\section{Experimental Setup}
\label{sec:setup}

\paragraph{Students and provenance labels.}
Every student is a low-rank (LoRA) fine-tune~\citep{hu2022lora} of the same base,
\textbf{Qwen2.5-7B-Instruct}~\citep{qwen2025qwen25}
(rank $32$, $\alpha\!=\!64$, dropout $0.05$ on all attention and MLP projections; $\sim\!1.05\%$ of
the $7.7$B parameters trainable), trained at sequence length $4096$ on
$N_{\mathrm{train}}=800$ problems with AdamW (learning rate $10^{-4}$, cosine schedule, $3\%$ warmup,
effective batch size $16$, bf16). The original CML grid uses the three-seed setting
reported below; the MTCR sensitivity sweep stores checkpoints at $E\in\{1,3,5\}$
distillation epochs, with $E=5$ serving as the full-duration setting. To avoid pseudoreplication we train each provenance type under
$3$ seeds (independent LoRA initialization, data order, and generation), giving $n\!=\!3$ student
draws per cell so detection is assessed at the \emph{student} level rather than by treating one
student's traces as independent. Because we generate every training corpus ourselves, provenance is
exact. The five provenance types are:
\begin{description}[leftmargin=1.4em,itemsep=1pt]
  \item[\texttt{suspect}] distilled from \textbf{DeepSeek-R1-Distill-Qwen-32B}~\citep{deepseekai2025r1} CoT traces (the
        protected teacher $T^\star$; a reasoning-style, cross-family teacher);
  \item[\texttt{ctrlInst}] distilled from \textbf{Qwen2.5-14B-Instruct} traces (a
        \emph{same-family}, non-reasoning teacher---the within-family case);
  \item[\texttt{ctrlLlama}] distilled from \textbf{DeepSeek-R1-Distill-Llama-8B} traces (a
        \emph{shared-lineage} R1 sibling on a different base---the teacher-specificity probe);
  \item[\texttt{ctrlC}] fine-tuned on \emph{human reference solutions} to the same problems, with
        \emph{no} teacher CoT (the quality-confound probe; non-distilled);
  \item[\texttt{ctrlA}] the base model itself, no fine-tuning (the clean non-distilled negative).
\end{description}

\paragraph{Tasks.}
The primary grid is \textbf{GSM8K}~\citep{cobbe2021gsm8k}: the train split supplies distillation/SFT problems and the
held-out test split the forensic trace pool ($N_{\mathrm{test}}=200$ problems, disjoint from SFT).
We run a single-seed stress test on \textbf{MATH}~\citep{hendrycks2021math}
(EleutherAI \texttt{hendrycks\_math}; longer, higher-entropy competition derivations,
sequence length $8192$) to check whether the trace-level pattern survives outside
GSM8K. The MTCR extension uses task-matched calibration views for sibling
attribution and the epoch-sensitivity sweep; the current submission reports these
views through aggregate source tables rather than a full per-task raw profile release.

\paragraph{Trace generation and scoring.}
Teacher and student traces are generated with vLLM~\citep{kwon2023vllm}
(bfloat16, temperature $0.6$, top-$p$ $0.95$).
Each student's $200$ held-out traces are scored under every candidate via the candidate's own
\texttt{prompt\_logprobs} over the assistant span, \emph{re-tokenized with the candidate's tokenizer}
(Eq.~\eqref{eq:scorelp}); this is the text-tier (T3) access regime. The surplus
(Eq.~\eqref{eq:surplus}) is taken relative to the base Qwen2.5-7B-Instruct as $b$.

\paragraph{Capability ladder.}
To probe the quality/coverage source of resemblance directly, we additionally score all students
under a \emph{same-family capability ladder}, Qwen2.5-Instruct at $\{1.5,3,7,14\}$B, none of which
generated any student's traces except via the $7$B base.

\paragraph{Evaluation.}
The independent replicate is the trained student. We report student-level summaries
for replicate separation and trace-level AUROC/TPR/FPR only as operating-curve
resolution within a fixed audit pool. Significance uses a $B=10{,}000$ label
permutation over the held-out trace scores, interpreted as an audit-operating test
rather than as additional independent training replicates. The null's KS-uniformity
check builds the null from one non-distilled control and evaluates the held-out
control's right-tail $p$-values against $\mathrm{Uniform}(0,1)$. Attribution is the
closed-set argmax over candidate teachers using the statistic being evaluated:
the base-relative surplus for the baseline, $r_c(S)$ for standard CML, and the
task-calibrated MTCR score for the MTCR attribution extension. We do not use this
argmax outside the declared candidate set.

\paragraph{Reviewer-stress result summaries.}
In addition to the raw GSM8K scored matrices, the revised manuscript imports locked
aggregate summaries for four reviewer stress families: adaptive trace transformations
(answer-only, compressed, paraphrased, mixed-source, and related variants across
GSM8K, MATH, BBH, ARC-Challenge, and MBPP), cross-base/cross-task transfer,
open-set decoys, and head-to-head baselines. These PI-provided summary tables
report adaptive AUROC, TPR@$1\%$FPR, FPR$_0$, confidence intervals, open-set
abstention and false-attribution rates, baseline metrics, and student-level
confidence intervals at the reported stratum level. The pasted tables do not carry
raw seed identities; we therefore treat them as aggregate result sources for
Figure~\ref{fig:pi_result_lock_suite} and the supplementary forensic-envelope
diagnostic.

\paragraph{Compute.}
The reported GSM8K focused grid (three teacher trace sets, four LoRA students, five
held-out trace sets, and scoring under six candidates) was run on a single MUSICA
H100 node and completes in
$\approx\!3.5$ minutes when generation and scoring are fanned across four GPUs. All models and
datasets are public. The source package contains the plotting scripts and aggregate
source tables needed to reproduce the manuscript figures. Code, plotting scripts,
aggregate source tables, and the manuscript source are available at
\url{https://github.com/Jie-Ni/CML-forensic-audits}.
